% !TEX program = xelatex
% Biên dịch bằng XeLaTeX, dùng cùng bộ phông với báo cáo BTL01.
\documentclass[11pt,a4paper]{article}

% ---------- Phông chữ (hỗ trợ đầy đủ tiếng Việt) ----------
\usepackage{fontspec}
\setmainfont{TeX Gyre Pagella}
\setsansfont{TeX Gyre Heros}
\setmonofont{Latin Modern Mono}[Scale=MatchLowercase]

% ---------- Trang giấy, tiêu đề ----------
\usepackage[a4paper,margin=2.5cm,headheight=22pt]{geometry}
\usepackage{microtype}
\usepackage{xcolor}
\usepackage{titlesec}
\usepackage{fancyhdr}
\usepackage{enumitem}
\usepackage[hidelinks]{hyperref}

% ---------- Bảng, công thức, mã ----------
\usepackage{array}
\usepackage{booktabs}
\usepackage{tabularx}
\usepackage{amsmath}
\usepackage{listings}

% ---------- Màu và kiểu chữ ----------
\definecolor{uplblue}{HTML}{1F3A5F}
\definecolor{uplgray}{HTML}{F3F5F8}
\definecolor{uplrule}{HTML}{C9D1DC}
\definecolor{upltok}{HTML}{8A2D1F}

\titleformat{\section}
  {\sffamily\Large\bfseries\color{uplblue}}
  {\thesection.}{0.6em}{}[\vspace{-0.4em}{\color{uplrule}\titlerule}]
\titlespacing*{\section}{0pt}{2.0ex plus 0.5ex}{1.2ex}

\pagestyle{fancy}
\fancyhf{}
\renewcommand{\headrulewidth}{0.3pt}
\renewcommand{\headrule}{\hbox to\headwidth{\color{uplrule}\leaders\hrule height \headrulewidth\hfill}}
\lhead{\small\sffamily\color{uplblue} Trình biên dịch UPL trên JVM}
\rhead{\small\sffamily\color{uplblue} BTL02}
\cfoot{\small\thepage}

\setlength{\parindent}{0pt}
\setlength{\parskip}{0.55em}
\renewcommand{\arraystretch}{1.25}
\setlist[itemize]{itemsep=0.25em,topsep=0.3em}
\setlist[enumerate]{itemsep=0.25em,topsep=0.3em}

\newcommand{\tm}[1]{\texttt{\textcolor{upltok}{#1}}}

% ---------- Mã nguồn UPL ----------
\lstdefinelanguage{UPL}{
  morekeywords={begin,end,int,bool,if,then,else,do,while,for,print,true,false},
  sensitive=true,
  morecomment=[l]{//},
  morecomment=[s]{/*}{*/},
}
\lstset{
  language=UPL,
  basicstyle=\ttfamily\small,
  keywordstyle=\color{uplblue},
  commentstyle=\itshape\color{gray},
  backgroundcolor=\color{uplgray},
  frame=single,
  rulecolor=\color{uplrule},
  framesep=6pt,
  xleftmargin=0pt,
  columns=fullflexible,
  keepspaces=true,
  showstringspaces=false,
  breaklines=true,
}

\begin{document}

\begin{flushleft}
  {\sffamily\bfseries\fontsize{22}{27}\selectfont\color{uplblue}
   TRÌNH BIÊN DỊCH UPL TRÊN JVM\\[2pt] PHÂN TÍCH NGỮ NGHĨA VÀ SINH MÃ\par}
  \vspace{0.8em}
  {\color{uplrule}\rule{\linewidth}{0.8pt}}\par
  \vspace{0.2em}
  \begin{tabular}{@{}l@{\hspace{0.6em}}l@{}}
    24021485 & Lã Việt Hoàng \\
    24021533 & Nguyễn Đức Khiêm \\
  \end{tabular}
\end{flushleft}

\section{Kiến trúc trình biên dịch}

Giai đoạn 2 sử dụng lại lexer Flex, parser Bison và AST của BTL01.
Trình biên dịch bổ sung pha kiểm tra ngữ nghĩa, biểu diễn mã ba địa chỉ
(Three-Address Code, TAC) và bộ sinh mã JVM theo phương án A của đề bài.

\[
  \begin{aligned}
    \text{UPL} &\xrightarrow{\text{Flex, Bison}} \text{AST}
    \xrightarrow{\text{ngữ nghĩa}} \text{AST có kiểu và ID biến} \\
    &\xrightarrow{\text{sinh TAC}} \text{TAC}
    \xrightarrow{\text{sinh Jasmin}} \texttt{.j}
    \xrightarrow{\text{Jasmin}} \texttt{.class}
  \end{aligned}
\]

\begin{tabularx}{\linewidth}{@{}l X@{}}
\toprule
\textbf{Tệp} & \textbf{Nhiệm vụ} \\
\midrule
\texttt{src/main.c} & Điều phối các pha và xử lý dòng lệnh \\
\texttt{src/semantic.c} & Bảng ký hiệu, phạm vi, liên kết định danh và kiểm tra kiểu \\
\texttt{src/tac.c} & Duyệt AST đã kiểm tra, sinh bộ bốn và xuất TAC dạng văn bản \\
\texttt{src/jvm.c} & Ánh xạ slot, ghi mã Jasmin và gọi trình lắp ráp \\
\bottomrule
\end{tabularx}

Chế độ \texttt{--compile} chạy toàn bộ các pha. Pha ngữ nghĩa bắt đầu khi lexer
và parser đã xử lý thành công. Khi phát hiện lỗi ngữ nghĩa, trình biên dịch báo
các lỗi tìm được và dừng trước khi sinh mã. Lớp đầu ra có phương thức
\texttt{public static void main(String[] args)} để chạy trên JVM.

\section{Bảng ký hiệu và phạm vi}

Mỗi mục trong \texttt{SymbolTable} lưu tên biến, kiểu \tm{int} hoặc \tm{bool},
ID phạm vi và dòng, cột của khai báo. Chỉ số của mục trong bảng là ID duy nhất
của biến. Chỉ số này được gắn vào node khai báo, gán hoặc đọc biến trong AST
và được dùng xuyên suốt khi sinh mã.

\begin{itemize}
  \item Chương trình có phạm vi gốc. Mỗi khối \texttt{\{ \}} mở phạm vi con.
        Nhánh \tm{then} và \tm{else} có phạm vi riêng. Điều kiện
        \tm{do-while} được kiểm tra trong phạm vi chứa vòng lặp.
  \item \tm{for} mở một phạm vi bao quanh phần khởi tạo, điều kiện, cập nhật
        và khối thân. Khai báo ở phần khởi tạo có hiệu lực đến hết vòng lặp.
        Khối thân vẫn có thể khai báo biến cùng tên trong phạm vi con.
  \item Khi dùng biến, trình biên dịch tìm từ phạm vi hiện tại ra các phạm vi
        cha và lấy khai báo gần nhất. Khi khai báo, chỉ kiểm tra trùng tên ở
        phạm vi hiện tại. Bảng vẫn giữ ID của biến sau khi rời phạm vi.
  \item Biểu thức khởi tạo được kiểm tra trước khi đưa biến mới vào bảng.
        Với \texttt{int x = x + 1;} trong khối con, \texttt{x} bên phải
        là biến ở phạm vi ngoài nếu biến đó đã được khai báo.
\end{itemize}

\clearpage
\section{Quy tắc ngữ nghĩa}

\begin{tabularx}{\linewidth}{@{}>{\raggedright\arraybackslash}p{4.6cm} X@{}}
\toprule
\textbf{Cấu trúc} & \textbf{Ràng buộc} \\
\midrule
Khai báo có khởi tạo, phép gán & Kiểu của biểu thức phải trùng kiểu biến \\
\addlinespace
Khai báo thiếu khởi tạo & Giá trị ban đầu là \tm{0} với \tm{int}, \tm{false} với \tm{bool} \\
\addlinespace
\tm{+}, \tm{*} & Hai toán hạng \tm{int}, kết quả \tm{int} \\
\tm{>}, \tm{>=} & Hai toán hạng \tm{int}, kết quả \tm{bool} \\
\tm{==} & Hai toán hạng cùng kiểu, kết quả \tm{bool} \\
\addlinespace
Điều kiện của \tm{if}, \tm{do-while}, \tm{for}
& Biểu thức điều kiện có kiểu \tm{bool} \\
\addlinespace
\tm{print} & Chấp nhận biểu thức hợp lệ kiểu \tm{int} hoặc \tm{bool} \\
\bottomrule
\end{tabularx}

Mỗi lỗi được in theo dạng \texttt{file:dòng:cột: semantic error: mô tả}.
Biểu thức lỗi được đánh dấu trong quá trình duyệt để tránh phát sinh chẩn
đoán sai kiểu dây chuyền. Hai toán hạng vẫn được kiểm tra độc lập, nên
\texttt{print(missing + other);} báo được cả hai tên chưa khai báo.

\section{Biểu diễn và sinh TAC}

TAC là dãy bộ bốn \texttt{(op, arg1, arg2, result)}. Toán hạng là hằng,
biến \texttt{vN}, biến tạm \texttt{tN}, nhãn \texttt{LN} hoặc trường rỗng.
Hằng, biến và biến tạm mang kiểu dữ liệu. Phần đầu tệp TAC liệt kê các biến
theo ID, kèm tên gốc, kiểu, phạm vi và vị trí khai báo.

\begin{tabularx}{\linewidth}{@{}>{\raggedright\arraybackslash}p{5.3cm} X@{}}
\toprule
\textbf{TAC} & \textbf{Cách sinh} \\
\midrule
\texttt{vN = x} & Ghi giá trị khởi tạo hoặc kết quả biểu thức gán vào biến \\
\texttt{tN = x op y} & Duyệt hai toán hạng trước, rồi ghi phép toán vào một biến tạm \\
\texttt{LN:}, \texttt{goto LN} & Đánh dấu đích hoặc chuyển điều khiển vô điều kiện \\
\texttt{if x goto LN} & Chuyển điều khiển khi điều kiện đúng \\
\texttt{ifFalse x goto LN} & Chuyển điều khiển khi điều kiện sai \\
\texttt{print x} & In giá trị của toán hạng theo kiểu \\
\bottomrule
\end{tabularx}

\begin{itemize}
  \item \tm{if} sinh \texttt{ifFalse} để bỏ qua nhánh \tm{then}. Nếu có
        \tm{else}, cuối nhánh \tm{then} sinh thêm \texttt{goto} đến nhãn
        hội tụ sau nhánh \tm{else}.
  \item \tm{do-while} đặt nhãn trước thân. Sau khi sinh thân và điều kiện,
        lệnh \texttt{if} quay lại đầu thân nếu điều kiện đúng.
  \item \tm{for} sinh khởi tạo trước nhãn kiểm tra. Sau điều kiện là
        \texttt{ifFalse} tới nhãn thoát. Khi điều kiện đúng, thân chạy trước
        cập nhật, rồi \texttt{goto} quay lại nhãn kiểm tra.
\end{itemize}

\clearpage
\section{Dịch TAC sang JVM}

Đoạn TAC mẫu của vòng \tm{do-while} dùng \texttt{v0} cho biến \texttt{sum}.

\begin{lstlisting}[language={}]
L4:
  t6:int = v0 + 1
  v0 = t6
  t7:bool = 26 > v0
  if t7 goto L4
\end{lstlisting}

Cả \tm{int} và \tm{bool} dùng nhóm lệnh nạp và lưu số nguyên của JVM.
Giá trị boolean được biểu diễn bằng \texttt{0} và \texttt{1}. Kiểu lưu trong
TAC quyết định hàm \texttt{println} được gọi.

\begin{tabularx}{\linewidth}{@{}>{\raggedright\arraybackslash}p{4cm} X@{}}
\toprule
\textbf{Dạng TAC} & \textbf{Lệnh JVM tương ứng} \\
\midrule
Hằng & \texttt{iconst\_n}, \texttt{bipush}, \texttt{sipush} hoặc
       \texttt{ldc\_w} tùy miền giá trị \\
Gán & Nạp hằng hoặc \texttt{iload}, sau đó \texttt{istore} \\
\tm{+}, \tm{*} & Nạp hai toán hạng, thực hiện \texttt{iadd} hoặc
                    \texttt{imul}, rồi lưu kết quả \\
\tm{>}, \tm{>=}, \tm{==} & \texttt{if\_icmpgt}, \texttt{if\_icmpge} hoặc
  \texttt{if\_icmpeq}, rồi nạp \texttt{0} hoặc \texttt{1} và lưu vào biến tạm \\
Nhãn và nhảy & Giữ nhãn TAC, dùng \texttt{goto\_w} để nhảy. Nhánh có điều
  kiện dùng \texttt{ifeq} hoặc \texttt{ifne} tới nhãn gần để bỏ qua
  \texttt{goto\_w} khi cần \\
\tm{print} & \texttt{getstatic} nạp \texttt{System.out}, nạp giá trị, rồi gọi
  \texttt{println(I)V} cho \tm{int} hoặc \texttt{println(Z)V} cho \tm{bool} \\
\bottomrule
\end{tabularx}

\subsection{Ánh xạ biến cục bộ}

Slot 0 chứa đối số \texttt{String[]} của \texttt{main}. Với $V$ biến UPL
và $T$ biến tạm, biến \texttt{v}$i$ dùng slot $1+i$ và biến tạm
\texttt{t}$j$ dùng slot $1+V+j$. Do đó \texttt{.limit locals} là $1+V+T$.
Ví dụ mẫu có $V=2$ và $T=10$, nên dùng 13 slot. Hai biến che khuất nhau có hai
ID khác nhau và luôn dùng hai slot khác nhau.

\subsection{Ngăn xếp toán hạng}

Ngăn xếp rỗng ở đầu và cuối mỗi bộ bốn. Phép toán hai ngôi nạp tối đa hai
giá trị, phép so sánh tiêu thụ hai toán hạng và hợp nhất hai nhánh nội bộ với
đúng một giá trị \texttt{0} hoặc \texttt{1} trước khi lưu. Lệnh in đặt đối
tượng \texttt{PrintStream} và giá trị trên ngăn xếp rồi tiêu thụ cả hai khi gọi
\texttt{println}. Do đó \texttt{.limit stack 2} đủ cho mọi bộ bốn.

Trình sinh mã kiểm tra giới hạn slot và số byte của phương thức \texttt{main}
trước khi lắp ráp. Các bước nhảy xa dùng \texttt{goto\_w} để xử lý thân
vòng lặp hoặc nhánh dài.

\clearpage
\section{Kiểm thử}

\texttt{make test} chạy bộ kiểm thử BTL01 và BTL02. Với chương trình hợp lệ,
bộ test lắp ráp, chạy \texttt{java -Xverify:all} và đối chiếu đầu ra với
tệp \texttt{.out}. Với lỗi ngữ nghĩa, bộ test so dòng, cột và nội dung với tệp
\texttt{.err}, đồng thời kiểm tra mã thoát và thư mục đầu ra.

\begin{tabularx}{\linewidth}{@{}>{\raggedright\arraybackslash}p{4.4cm} X X@{}}
\toprule
\textbf{Ca kiểm thử} & \textbf{Kết quả mong đợi} & \textbf{Kết quả thực tế} \\
\midrule
\texttt{01\_control\_flow} & In \texttt{26}, \texttt{true} & Khớp cả hai dòng \\
\texttt{02\_scopes} & In giá trị của biến theo phạm vi và giá trị mặc định
  & Khớp 9 dòng trong \texttt{.out} \\
\texttt{03\_expressions} & Đúng độ ưu tiên, so sánh, giá trị biên và số lượt lặp
  & Khớp 11 dòng trong \texttt{.out} \\
\texttt{02\_empty} & JVM chạy thành công, đầu ra rỗng & Mã thoát 0, đầu ra rỗng \\
\texttt{01\_scopes} (lỗi) & 8 lỗi liên kết tên và phạm vi
  & Khớp 8 chẩn đoán, không sinh mã \\
\texttt{02\_types} (lỗi) & 14 lỗi kiểu và định danh
  & Khớp 14 chẩn đoán, không sinh mã \\
Lỗi lexer và parser khi biên dịch & Mã thoát 1, dừng trước pha ngữ nghĩa
  & Cả hai ca đều đạt \\
\bottomrule
\end{tabularx}

Bộ kiểm thử đạt 33/33 với Jasmin 2.5.0 và OpenJDK 21. Các ca kiểm tra thêm
xác nhận lớp dùng hơn 255 slot và có bước nhảy xa chạy được trên JVM.
Trình biên dịch từ chối phương thức \texttt{main} vượt giới hạn 65535 byte.

Các tệp \texttt{UplProgram.tac}, \texttt{UplProgram.j} và
\texttt{UplProgram.class} trong \texttt{examples/phase2/} được sinh từ
\texttt{01\_control\_flow.upl}.

\begin{lstlisting}[language={}]
make
build/upl --compile build/btl02-demo tests/phase2/valid/01_control_flow.upl
java -Xverify:all -cp build/btl02-demo UplProgram
\end{lstlisting}

Kết quả thực tế của lớp mẫu gồm hai dòng.

\begin{lstlisting}[language={}]
26
true
\end{lstlisting}

\end{document}
